\documentclass[11pt]{article}
\usepackage[margin=0.9in]{geometry}
\usepackage{amsmath,amssymb,amsthm,mathtools}
\usepackage[T1]{fontenc}
\usepackage[utf8]{inputenc}
\usepackage{enumitem}
\usepackage{booktabs,longtable,array,seqsplit}
\usepackage{hyperref}
\usepackage{xcolor}
\usepackage{microtype}
\hypersetup{colorlinks=true,linkcolor=blue!45!black,citecolor=blue!45!black,urlcolor=blue!45!black}
\setlength{\emergencystretch}{6em}

\newtheorem{theorem}{Theorem}[section]
\newtheorem{proposition}[theorem]{Proposition}
\newtheorem{lemma}[theorem]{Lemma}
\newtheorem{corollary}[theorem]{Corollary}
\newtheorem{claim}[theorem]{Claim}
\theoremstyle{definition}
\newtheorem{definition}[theorem]{Definition}
\newtheorem{assumption}[theorem]{Assumption}
\theoremstyle{remark}
\newtheorem{remark}[theorem]{Remark}

\newcommand{\R}{\mathbb{R}}
\newcommand{\C}{\mathbb{C}}
\newcommand{\N}{\mathbb{N}}
\newcommand{\Hsp}{\mathcal{H}}
\newcommand{\Fsp}{\mathcal{F}}
\newcommand{\Id}{\mathrm{Id}}
\newcommand{\tr}{\operatorname{tr}}
\newcommand{\diag}{\operatorname{diag}}
\newcommand{\rank}{\operatorname{rank}}
\newcommand{\spec}{\operatorname{spec}}
\newcommand{\spanop}{\operatorname{span}}
\newcommand{\dist}{\operatorname{dist}}
\newcommand{\range}{\operatorname{ran}}
\newcommand{\Ker}{\operatorname{ker}}
\newcommand{\Pinch}{\mathcal{C}}
\newcommand{\Off}{\widetilde{\mathcal{C}}}
\newcommand{\Pperp}{P^{\perp}}
\newcommand{\Qperp}{Q^{\perp}}
\newcommand{\op}{\mathrm{op}}
\newcommand{\F}{\mathrm{F}}
\newcommand{\HS}{\mathrm{HS}}
\newcommand{\ip}[2]{\left\langle #1,#2\right\rangle}
\newcommand{\norm}[1]{\left\lVert #1\right\rVert}
\newcommand{\abs}[1]{\left\lvert #1\right\rvert}
\newcommand{\code}[1]{\nolinkurl{#1}}

\newenvironment{argument}{\paragraph{Argument.}\begin{enumerate}[leftmargin=2em]}{\end{enumerate}}
\newenvironment{proofoutline}{\paragraph{Proof architecture.}\begin{enumerate}[leftmargin=2em]}{\end{enumerate}}
\newcommand{\sourceanchor}[1]{\par\smallskip\noindent\textbf{Source anchor.} #1\par\smallskip}
\newcommand{\sourcecomment}[1]{\begin{remark}#1\end{remark}}
\newenvironment{sourceguide}{\begin{description}[style=nextline,leftmargin=3.3cm,labelwidth=3.1cm]}{\end{description}}
\newenvironment{formalizationmap}{\begin{longtable}{@{}p{0.20\textwidth}p{0.31\textwidth}p{0.40\textwidth}@{}}\toprule Source item & Modern mathematical content & Repository status \\\midrule\endfirsthead\toprule Source item & Modern mathematical content & Repository status \\\midrule\endhead}{\bottomrule\end{longtable}}

\title{Grubi\v{s}i\'c--Kostrykin--Makarov--Veseli\'c (2013): tan $2\Theta$ through signs, polar factors, and indefinite forms}
\author{}
\date{Distilled reconstruction for the finite Davis--Kahan formalization}
\begin{document}
\maketitle

\section*{Source map and scope}
\begin{sourceguide}
\item[Primary source] Luka Grubi\v{s}i\'c, Vadim Kostrykin, Konstantin A. Makarov, and Kre\v{s}imir Veseli\'c, \emph{The Tan $2\Theta$ Theorem for Indefinite Quadratic Forms}, Journal of Spectral Theory 3 (2013), 83--100, DOI \href{https://doi.org/10.4171/JST/38}{10.4171/JST/38}; preprint \href{https://arxiv.org/abs/1006.3190}{arXiv:1006.3190}.
\item[Paper scope] Possibly unbounded, not necessarily semibounded self-adjoint operators defined by closed quadratic forms; off-diagonal form perturbations; sign operators and positive spectral projections; sharp quarter-turn estimates.
\item[Formalization target] The bounded finite specialization underlying \code{TanTwoTheta.lean}: off-diagonal perturbations, reflection anticommutation, sign/projector geometry, the sharp factor $2$, and proof of \code{AvoidsQuarterTurn}.
\item[Source-faithful norm scope] The principal theorem is an operator-norm/maximal-angle result. The repository may derive stronger finite UI-norm descendants by separate finite majorization arguments; those strengthenings should not be retroactively attributed to this paper.
\end{sourceguide}

The paper reveals a conceptual form of tan $2\Theta$ that is especially useful for formalization. Instead of beginning with a tangent operator, it compares two self-adjoint involutions: the sign of the unperturbed operator and the sign of the perturbed operator. The angle bound follows from the polar factor of a sectorial operator. This simultaneously proves the numerical estimate and the strict quarter-turn branch.

\section{The bounded classical model}

Let
\[
 A=
 \begin{pmatrix}
 A_+&0\\
 0&-A_-
 \end{pmatrix}
\]
on $\Hsp=\Hsp_+\oplus\Hsp_-$, with $A_+>0$ and $A_->0$. Let
\[
 J=
 \begin{pmatrix}
 I&0\\
 0&-I
 \end{pmatrix}
\]
be the grading involution. An off-diagonal self-adjoint perturbation has the form
\[
 V=
 \begin{pmatrix}
 0&B\\
 B^*&0
 \end{pmatrix},
\]
equivalently
\[
 JV=-VJ.
\]
Set $L=A+V$. The reference projection and perturbed positive spectral projection are
\[
 P=E_A((0,\infty))=\frac{I+J}{2},
 \qquad
 Q=E_L((0,\infty))=\frac{I+\operatorname{sign} L}{2}.
\]

The classical tan $2\Theta$ estimate reads
\[
 \norm{\tan(2\Theta)}_{\op}
 \le\frac{2\norm V_{\op}}{d},
\]
where $d$ is the spectral separation between the two diagonal blocks in the appropriate shifted configuration. Equivalently,
\[
 \norm{P-Q}_{\op}
 \le
 \sin\left(\frac12\arctan\frac{2\norm V}{d}\right)
 <\frac{1}{\sqrt2}.
\]
The strict final inequality is not an afterthought. It says the perturbed positive subspace lies on the same side of the quarter-turn branch as the reference subspace.

\section{Quadratic-form formulation}

The paper replaces bounded operators by a symmetric quadratic form $\mathfrak a$ and a self-adjoint involution $J$. Its basic hypothesis requires
\[
 \mathfrak a_J[x,y]:=\mathfrak a[x,Jy]
\]
to be positive definite and closed. The eigenspaces of $J$ supply the positive and negative reference decomposition.

A symmetric perturbation form $\mathfrak v$ is off-diagonal when
\[
 \mathfrak v[Jx,y]=-\mathfrak v[x,Jy].
\]
The relative size is measured by
\[
 v_0
 =
 \sup_{x\ne0}
 \frac{\abs{\mathfrak v[x]}}{\mathfrak a_J[x]}.
\]
For a shift $\mu$ in the unperturbed gap, define
\[
 \mathfrak a_\mu[x,y]
 =\mathfrak a[x,Jy]-\mu\ip{x}{Jy},
 \qquad
 v_\mu
 =
 \sup_{x\ne0}
 \frac{\abs{\mathfrak v[x]}}{\mathfrak a_\mu[x]}.
\]
The optimization over $\mu$ is the form-theoretic analogue of choosing the best center of a finite spectral gap.

In finite dimensions, all forms are bounded and represented by matrices. Nevertheless this formulation explains which assumptions are structural:
\begin{itemize}[leftmargin=2em]
\item a grading/reflection $J$;
\item positivity of the signed unperturbed operator $JA$ after a shift;
\item off-diagonal anticommutation of the perturbation;
\item a relative perturbation size.
\end{itemize}

\section{The polar-factor mechanism}

For each admissible shift $\mu$, the paper constructs a sectorial operator $T_\mu$ representing the perturbed form relative to the positive form $\mathfrak a_\mu$. Its polar decomposition is
\[
 T_\mu=U_\mu\abs{T_\mu}.
\]
The sectorial semi-angle is bounded by
\[
 \theta_\mu=\arctan v_\mu.
\]
A key structural identity connects the perturbed self-adjoint operator $L-\mu$ to $T_\mu$:
\[
 L-\mu=JT_\mu
\]
at the form/representation level. Consequently
\[
 \abs{T_\mu}=\abs{L-\mu},
 \qquad
 U_\mu=J\operatorname{sign}(L-\mu).
\]
Because $\mu$ remains inside a resolvent gap, $\operatorname{sign}(L-\mu)$ is independent of the particular shift.

The polar factor of a sectorial operator stays in the same sector. Therefore
\[
 \norm{I-U_\mu}
 \le
 2\sin\left(\frac12\arctan v_\mu\right).
\]
Using $U_\mu=J\operatorname{sign} L$ gives
\[
 \norm{J-\operatorname{sign} L}
 \le
 2\sin\left(\frac12\arctan v_\mu\right).
\]
Since
\[
 P-Q=\frac{J-\operatorname{sign} L}{2},
\]
one obtains the projector estimate
\begin{equation}
 \norm{P-Q}
 \le
 \sin\left(\frac12\arctan v\right)
 <\frac1{\sqrt2},
 \qquad
 v:=\inf_\mu v_\mu.
 \tag{GKMMV-proj}\label{GKMMV-proj}
\end{equation}

This proof packages branch selection automatically: $\arctan v<\pi/2$, so the half-angle is strictly below $\pi/4$.

\section{Theorem 3.1: indefinite sign-subspace rotation}

Let $A$ and $L$ be the self-adjoint operators associated with the unperturbed and perturbed forms. Let
\[
 P=E_A((0,\infty)),
 \qquad
 Q=E_L((0,\infty)).
\]
Under the form hypotheses, Theorem 3.1 states exactly the bound \eqref{GKMMV-proj}. In angle language,
\[
 \norm\Theta(P,Q)
 \le
 \frac12\arctan v
 <\frac{\pi}{4}.
\]
Equivalently,
\[
 \norm{\tan(2\Theta)}\le v.
\]
For the bounded additive model, $v$ reduces to a normalized off-diagonal perturbation size and recovers the familiar factor $2\norm V/d$.

\section{Semibounded operators with a finite gap}

The paper next treats a semibounded self-adjoint operator with a finite gap $(\alpha,\beta)$. Let the reference spectral projection select one side of the gap, and assume the perturbation form is off-diagonal relative to this projection.

With distances $d_+$ and $d_-$ from the chosen shift/selected spectrum to the two gap edges, Theorem 4.2 gives
\[
 \norm{P-Q}
 \le
 \sin\left(
   \frac12\arctan\frac{2v}{\delta}
 \right)
 <\frac1{\sqrt2},
\]
where the dimensionless geometric factor $\delta$ is built from $d_+$ and $d_-$; the paper distinguishes two configurations, producing respectively
\[
 \delta=\frac{d_++d_-}{\sqrt{d_+d_-}}
 \quad\text{or}\quad
 \delta=\frac{d_+-d_-}{\sqrt{d_+d_-}}.
\]
The first configuration is the stable interior-gap analogue of the classical theorem; the second records a more delicate asymmetric arrangement.

A two-dimensional example attains the displayed angle. This confirms both the factor $2$ and the half-arctangent functional form. Any finite theorem that changes the constant or drops the branch conclusion must therefore be clearly labeled as a different result.

\section{Reflection and sign dictionaries for Lean}

The finite formalization can avoid quadratic-form infrastructure while retaining the proof's reusable core.

Given a subspace $U$, define its reflection
\[
 J_U=2P_U-I.
\]
For a self-adjoint operator $L$ with no zero eigenvalue,
\[
 \operatorname{sign} L=2E_L((0,\infty))-I.
\]
Thus the difference of projections and the difference of involutions satisfy
\[
 2(P_U-P_V)=J_U-J_V.
\]
The direct-rotation phase is encoded by
\[
 J_UJ_V,
\]
whose eigenvalues on a principal two-plane are $e^{\pm2i\theta}$. This is the exact origin of the doubled angle.

The following equivalences are high-value API lemmas:
\begin{align*}
 \norm{P_U-P_V}<1/\sqrt2
 &\Longleftrightarrow
 \theta_{\max}<\pi/4,\\
 \theta_{\max}<\pi/4
 &\Longrightarrow
 \cos(2\Theta)>0,\\
 \cos(2\Theta)>0
 &\Longrightarrow
 \tan(2\Theta)\text{ is everywhere defined}.
\end{align*}
They separate the branch proof from the subsequent norm estimate.

\section{Relation to the Riccati proof}

When $V$ is off-diagonal and $Q$ is a graph subspace, its graph operator $X$ solves a Riccati equation. The tangent identities are
\[
 \tan\Theta=\abs X,
 \qquad
 \tan(2\Theta)=2\abs X(I-\abs X^2)^{-1}.
\]
A Riccati manipulation yields a Sylvester equation whose unknown is the doubled tangent expression. The paper's sign/polar proof and the finite Riccati/Sylvester proof are two views of the same geometry:
\begin{itemize}[leftmargin=2em]
\item the sign proof establishes the acute branch globally;
\item the Riccati proof exposes the residual and is well suited to UI-norm estimates;
\item both use off-diagonality essentially.
\end{itemize}

For the current Lean development, it may be best to prove branch selection with reflection/sign or spectral repulsion, then prove the quantitative tangent estimate with the existing finite Sylvester/Ky Fan machinery.

\section{Finite formalization map}

\begin{formalizationmap}
Grading & $J=2P-I$, a self-adjoint isometric involution & Projection/reflection lemmas should be centralized.\\
Off-diagonal perturbation & $JV=-VJ$, equivalently diagonal compressions vanish & Existing hypotheses around \code{offDiagonal} and residual identities should expose this equivalence.\\
Sign projection & $E_L((0,\infty))=(I+\operatorname{sign} L)/2$ & Finite spectral-projector wrappers can express this through eigenvalue selection.\\
Quarter-turn avoidance & $\norm{P-Q}<1/\sqrt2$ implies every angle is $<\pi/4$ & \code{isAcute_canonical_tanTwoTheta} and \code{largestPrincipalAngle_lt_pi_div_four} target this content.\\
Raw tan $2\Theta$ & $\delta N(\tan2\Theta)\le2N(V)$ & \code{tanTwoTheta_perturbation_le} and norm-specific wrappers are scaffolded.\\
Canonical spectral theorem & Choose the perturbed spectral subspace and prove branch/gap persistence & \code{tanTwoTheta_spectralSubspace_le} is the intended wrapper.\\
Sharpness & Two-dimensional off-diagonal model attains half-arctangent bound & Should become a finite regression example after the endpoint is complete.\\
\end{formalizationmap}

\section{Proof guidance for the remaining finite endpoint}

A stable route is:
\begin{enumerate}[leftmargin=2em]
\item Represent the reference subspace by its reflection $J$ and show the perturbation anticommutes with $J$.
\item Choose a shift in the gap so the signed diagonal part $J(A-\mu I)$ is positive.
\item Prove spectral repulsion or a positive lower bound for the perturbed signed operator.
\item Deduce $\norm{P-Q}<1/\sqrt2$ and hence \code{AvoidsQuarterTurn}.
\item Construct the graph/angular operator only after the branch has been established.
\item Derive the Riccati identity and isolate the linear Sylvester equation for the doubled tangent embedding.
\item Apply the ordered-gap Ky Fan/UI estimate, yielding the factor $2$.
\item Add the operator-norm Part III alias and the strict quarter-turn conclusion together.
\end{enumerate}

\section{Failure modes to exclude}

\begin{itemize}[leftmargin=2em]
\item A bound on $\sin(2\Theta)$ alone does not choose between angles below and above $\pi/4$; branch selection needs a separate argument.
\item Off-diagonal means anticommutation with the reference reflection, not merely that one displayed matrix block happens to vanish in an arbitrary decomposition.
\item Do not divide by $\cos(2\Theta)$ before proving it is invertible and positive on the selected branch.
\item Do not attribute an all-UI-norm finite strengthening to this operator-norm form theorem.
\item Do not discard the factor $2$: it is intrinsic to the double-angle identity and is sharp in dimension two.
\end{itemize}

\section*{Distilled conclusion}

The most reusable lesson of the paper is that tan $2\Theta$ is fundamentally a theorem about two sign involutions. Off-diagonal perturbation makes the perturbed sign accessible through a sectorial polar factor; the polar angle gives both the sharp quantitative estimate and strict avoidance of the quarter turn. In the finite Lean theory, this sign/reflection layer should establish the branch, while Riccati and Sylvester machinery supplies the source-faithful residual and norm inequalities.

\end{document}
